\chapter{PDF condizionate e funzioni di variabili continue}

Questo capitolo porta nel continuo due strumenti che conosciamo già dal discreto (Capitolo~6). Il primo è il
\textbf{condizionamento}: come cambia la densità di $X$ quando si sa che si è verificato un evento. Il secondo sono le
\textbf{funzioni di variabili aleatorie}: se $X$ ha una certa densità, che densità ha $Y = g(X)$? La risposta è una
formula che nel discreto non c'era, perché trasformare una variabile continua non sposta soltanto la probabilità, ma la
\emph{stira} e la \emph{comprime}. Vedremo anche due applicazioni importantissime: come un computer genera numeri
casuali con qualsiasi densità, e come un convertitore analogico-digitale trasforma una variabile continua in una
discreta.

% ============================================================
\section{PDF e CDF condizionate a un evento}
% ============================================================

In modo del tutto analogo a quanto fatto per le variabili discrete, se $A$ è un evento con $P(A) > 0$ possiamo
considerare la probabilità che $X$ cada in un intervallino attorno a $x$ \emph{sapendo} che si è verificato $A$:
\[ P\left( X \in \left[ x - \tfrac{\Delta x}{2},\ x + \tfrac{\Delta x}{2} \right] \;\Big|\; A \right) = P_X(x, \Delta x
   \mid A) . \]

\dfn{PDF e CDF condizionate}{
  La \term{densità di $X$ condizionata all'evento $A$} è
  \[ f_{X \mid A}(x) = \lim_{\Delta x \to 0} \frac{P_X(x, \Delta x \mid A)}{\Delta x} . \]
  La \term{CDF condizionata} è
  \[ F_{X \mid A}(x) = P(X \leq x \mid A) = \frac{P(\{X \leq x\} \cap A)}{P(A)} , \qquad\text{e}\qquad
     f_{X \mid A}(x) = \frac{\dd}{\dd x} F_{X \mid A}(x) . \]
}

Siccome la probabilità condizionata è a sua volta una probabilità (Teorema~\ref{teo:condprob}), $f_{X \mid A}$ è una
densità a tutti gli effetti: non negativa e con area 1. Il caso più frequente è quello in cui $A$ dice che $X$ è caduta
in un certo intervallo.

\begin{esempio}{Una Laplace ristretta a un intervallo}
Sia $X \sim \mathcal{L}(\lambda)$ e sia $A = \{-1 \leq X \leq 2\}$. Allora
\[ F_{X \mid A}(x) = P(X \leq x \mid A) = \frac{P(\{X \leq x\} \cap \{-1 \leq X \leq 2\})}{P(-1 \leq X \leq 2)} , \]
e il denominatore è $P(A) = F_X(2) - F_X(-1)$. Per il numeratore si guarda dove si trova $x$:
\begin{itemize}
  \item se $x < -1$, i due eventi sono incompatibili (non si può essere sotto $-1$ e fra $-1$ e 2): l'intersezione è
        vuota;
  \item se $-1 \leq x < 2$, l'intersezione è $\{-1 \leq X \leq x\}$, con probabilità $F_X(x) - F_X(-1)$;
  \item se $x \geq 2$, l'intersezione è tutto $A$.
\end{itemize}
Quindi
\[
  F_{X \mid A}(x) =
  \begin{cases}
    0 & x < -1 \\[1mm]
    \dfrac{F_X(x) - F_X(-1)}{F_X(2) - F_X(-1)} & -1 \leq x < 2 \\[3mm]
    \dfrac{F_X(2) - F_X(-1)}{F_X(2) - F_X(-1)} = 1 & x \geq 2
  \end{cases}
  \qquad\qquad
  f_{X \mid A}(x) =
  \begin{cases}
    \dfrac{f_X(x)}{F_X(2) - F_X(-1)} & -1 \leq x < 2 \\[3mm]
    0 & \text{altrove}
  \end{cases}
\]
La densità condizionata si ottiene derivando: il denominatore è una costante. Con $\lambda = 1$:
$F_X(2) = 1 - \frac12 e^{-2} \approx 0{,}932$, $F_X(-1) = \frac12 e^{-1} \approx 0{,}184$, e $P(A) \approx 0{,}748$.
\begin{center}
\begin{tikzpicture}
\begin{axis}[quaderno, width=11cm, height=5cm, xmin=-3.6, xmax=3.8, ymin=0, ymax=0.75, xtick={-3,-2,-1,1,2,3},
  ytick={0.5,0.668}, yticklabels={$0{,}5$,$0{,}67$}, xlabel=$x$, legend style={font=\scriptsize, at={(0.98,0.98)},
  anchor=north east}]
  \addplot[fill=evid, draw=none, domain=-1:2, samples=60, forget plot] {0.5*exp(-abs(x))} \closedcycle;
  \addplot[mspiA, domain=-3.4:3.4, samples=161] {0.5*exp(-abs(x))};
  \addplot[mspiB, line width=1.3pt, domain=-1:2, samples=61] {0.5*exp(-abs(x))/0.7484};
  \addplot[mspiB, line width=1.3pt, forget plot] coordinates {(-3.4,0) (-1,0)};
  \addplot[mspiB, line width=1.3pt, forget plot] coordinates {(2,0) (3.4,0)};
  \draw[mspiB, dashed] (axis cs:-1,0) -- (axis cs:-1,0.2459);
  \draw[mspiB, dashed] (axis cs:2,0) -- (axis cs:2,0.0904);
  \legend{$f_X(x)$, $f_{X \mid A}(x)$}
\end{axis}
\end{tikzpicture}
\end{center}
L'area gialla vale $P(A) \approx 0{,}748$. La densità condizionata (in arancione) è la stessa curva, tagliata fuori
dall'intervallo e «rialzata» dividendola per $0{,}748$: così la sua area torna a essere 1. È esattamente quello che
succedeva con le pmf nel Capitolo~6.
\end{esempio}

% ============================================================
\section{La probabilità totale per le densità}
% ============================================================

\teorema[totalepdf]{Probabilità totale per le PDF e per le medie}{
  Se $\{E_m\}_{m=1}^{M}$ è una partizione di $\Omega$ con $P(E_m) > 0$, allora
  \[ F_X(x) = \sum_{m=1}^{M} F_{X \mid E_m}(x)\, P(E_m) , \qquad f_X(x) = \sum_{m=1}^{M} f_{X \mid E_m}(x)\, P(E_m) , \]
  \[ E[X] = \sum_{m=1}^{M} E[X \mid E_m]\, P(E_m) = \sum_{m=1}^{M} P(E_m) \int_{-\infty}^{\infty} x\, f_{X \mid E_m}(x) \dd x . \]
}

\begin{dimostrazione}
\passo{La CDF} Come nel Capitolo~3, si ritaglia l'evento con la partizione:
\[
  F_X(x) = P(\{X \leq x\} \cap \Omega) = P\left( \{X \leq x\} \cap \bigcup_{m=1}^{M} E_m \right)
  = P\left( \bigcup_{m=1}^{M} \left( \{X \leq x\} \cap E_m \right) \right) ,
\]
e i pezzi sono incompatibili, quindi le probabilità si sommano:
\[ F_X(x) = \sum_{m=1}^{M} P(\{X \leq x\} \cap E_m) = \sum_{m=1}^{M} F_{X \mid E_m}(x)\, P(E_m) . \]
\passo{La densità} Si deriva rispetto a $x$: la derivata di una somma è la somma delle derivate, e i $P(E_m)$ sono
costanti.
\passo{La media} Si moltiplica la densità per $x$ e si integra: per la linearità dell'integrale
\[ E[X] = \int x \sum_m f_{X \mid E_m}(x)\, P(E_m) \dd x = \sum_m P(E_m) \int x\, f_{X \mid E_m}(x) \dd x = \sum_m
   P(E_m)\, E[X \mid E_m] . \]
\end{dimostrazione}

\begin{esempio}{La Laplace spezzata in due}
Con riferimento all'esempio precedente ($\lambda = 1$), gli eventi $A$ e $\compl{A}$ formano una partizione:
\[ E[X] = E[X \mid A]\, P(A) + E[X \mid \compl{A}]\, P(\compl{A}) . \]
Ogni pezzo ha un significato semplice:
$E[X \mid A]\, P(A) = \frac{1}{P(A)} \int_{-1}^{2} x\, f_X(x) \dd x \cdot P(A) = \int_{-1}^{2} x\, f_X(x) \dd x$, cioè
l'integrale di $x f_X(x)$ ristretto ai valori per cui $A$ è vero; e lo stesso per $\compl{A}$, con i valori per cui $A$ è
falso. Sommando i due pezzi si integra su tutta la retta, e si ritrova $E[X] = \int x\, f_X(x) \dd x = 0$.

Con i numeri: usando $\int_0^c x\, e^{-x} \dd x = 1 - (1 + c)\, e^{-c}$ (per parti),
\[ \int_{-1}^{2} x\, \tfrac12 e^{-\abs{x}} \dd x = \tfrac12 \left[ \int_0^2 x\, e^{-x} \dd x - \int_0^1 x\, e^{-x} \dd x
   \right] = \tfrac12 \left[ (1 - 3e^{-2}) - (1 - 2e^{-1}) \right] \approx 0{,}165 , \]
quindi $E[X \mid A] \approx \frac{0{,}165}{0{,}748} \approx 0{,}22$: sapendo che $X$ è fra $-1$ e 2, la media si sposta
verso destra, perché l'intervallo è più lungo da quella parte. Il pezzo di $\compl{A}$ vale $-0{,}165$, e il totale è
zero.
\end{esempio}

\begin{esempio}{Due tipi di componenti}
Un magazzino contiene componenti di due fornitori, in parti uguali. Quelli del primo durano un tempo esponenziale di
media 1 anno ($\lambda = 1$), quelli del secondo di media 3 anni ($\lambda = \frac13$). Si prende un componente a caso:
la sua durata ha densità
\[ f_X(x) = \tfrac12 \cdot 1 \cdot e^{-x} u(x) + \tfrac12 \cdot \tfrac13 e^{-x/3} u(x) , \]
una «miscela» di due esponenziali, e media $E[X] = \frac12 \cdot 1 + \frac12 \cdot 3 = 2$ anni, senza calcolare nessun
integrale.
\end{esempio}

% ============================================================
\section{Funzioni di variabili aleatorie continue}
% ============================================================

Il problema è simile a quello già affrontato per le variabili discrete. Si assuma che $X = X(\omega)$ sia una variabile
aleatoria continua con alfabeto $\mathcal{X}$, densità $f_X(x)$ e CDF $F_X(x)$, e sia $g$ una funzione definita sui
punti di $\mathcal{X}$ (tranne, al più, un insieme di probabilità nulla). Si forma la nuova variabile aleatoria
\[ Y = g(X) = g[X(\omega)] \in \mathcal{Y} , \qquad \mathcal{Y} = g(\mathcal{X}) . \]
Vogliamo ricavare la caratterizzazione di $Y$ da quella di $X$: la densità (o la CDF) e, nel prossimo capitolo, la
media. A differenza del caso discreto, qui si distinguono tre casi:
\begin{enumerate}[label=\alph*)]
  \item $g$ è \textbf{biunivoca}, cioè invertibile, continua e derivabile;
  \item $g$ non è invertibile, ma è continua e derivabile, e $\mathcal{Y}$ è continuo;
  \item $g$ non è invertibile e $\mathcal{Y}$ è \textbf{discreto}: è una \textbf{conversione A/D} della variabile
        continua (una quantizzazione, cioè una compressione con perdite), come quella che si fa sui segnali
        analogici.
\end{enumerate}

\begin{richiamo}{funzioni monotone e derivata della funzione inversa}
Una funzione è \term{strettamente crescente} se a $x$ più grandi corrispondono valori più grandi, \term{strettamente
decrescente} se corrispondono valori più piccoli; in entrambi i casi si dice \term{strettamente monotona}. Una funzione
continua su un intervallo è invertibile esattamente quando è strettamente monotona: se tornasse indietro, qualche
valore verrebbe raggiunto due volte. L'inversa $x = g^{-1}(y)$ risponde alla domanda «quale $x$ dà $y$?», e il suo
grafico è quello di $g$ riflesso rispetto alla bisettrice $y = x$.

\textbf{Derivata dell'inversa.} Per definizione $g\big(g^{-1}(y)\big) = y$. Derivando entrambi i membri rispetto a $y$
con la regola della catena: $\dot g\big(g^{-1}(y)\big) \cdot \frac{\dd}{\dd y} g^{-1}(y) = 1$, cioè
\[ \dot x(y) = \frac{\dd}{\dd y}\, g^{-1}(y) = \frac{1}{\dot g\big(g^{-1}(y)\big)} : \]
la derivata della funzione inversa in un punto è l'inverso della derivata della funzione originaria, valutata nel punto
corrispondente. (Il puntino, $\dot g$, è un altro modo di indicare la derivata, usato nelle slide.) Nel disegno: le due
tangenti sono una il riflesso dell'altra, e scambiare $x$ con $y$ trasforma la pendenza $m$ in $\frac1m$.
\begin{center}
\begin{tikzpicture}
\begin{axis}[quaderno, width=6.4cm, height=6.4cm, xmin=0, xmax=2.6, ymin=0, ymax=2.6, xtick=\empty, ytick=\empty,
  grid=none, axis equal image]
  \addplot[mspigray, dashed, domain=0:2.5] {x};
  \addplot[mspiA, domain=0:1.58, samples=40] {x^2};
  \addplot[mspiB, domain=0:2.5, samples=60] {sqrt(x)};
  \addplot[mspiA!85!black, domain=0.8:1.6] {1.44+2.4*(x-1.2)};
  \addplot[mspiB!85!black, domain=0.6:2.3] {1.2+(x-1.44)/2.4};
  \fill[mspiA] (axis cs:1.2,1.44) circle (1.6pt); \fill[mspiB] (axis cs:1.44,1.2) circle (1.6pt);
  \node[font=\scriptsize, mspiA!85!black, anchor=east] at (axis cs:0.95,1.95) {$g(x) = x^2$};
  \node[font=\scriptsize, mspiB!85!black, anchor=north west] at (axis cs:1.6,1.0) {$g^{-1}(y) = \sqrt{y}$};
  \node[font=\scriptsize, mspigray, anchor=south] at (axis cs:2.3,2.35) {$y = x$};
\end{axis}
\end{tikzpicture}
\hspace{0.5cm}
\begin{minipage}[b]{6.6cm}\small
In $x = 1{,}2$ la parabola ha pendenza $2x = 2{,}4$. Nel punto corrispondente $y = 1{,}44$ la radice ha pendenza
$\frac{1}{2\sqrt{y}} = \frac{1}{2{,}4}$. Le pendenze sono una l'inversa dell'altra.
\end{minipage}
\end{center}
\end{richiamo}

\subsection{Primo caso: \texorpdfstring{$g$}{g} è invertibile}

Se $g$ è invertibile è strettamente monotona su $\mathcal{X}$.

\teorema[monotona]{Densità di una funzione monotona}{
  Se $g$ è strettamente monotona e derivabile, la densità di $Y = g(X)$ è
  \[ f_Y(y) = \frac{f_X\big(g^{-1}(y)\big)}{\abs{\dot g\big(g^{-1}(y)\big)}} = \frac{f_X\big(x(y)\big)}{\abs{\dot
     g\big(x(y)\big)}} , \]
  dove $x(y) = g^{-1}(y)$ è la soluzione dell'equazione $y = g(x)$, che naturalmente dipende da $y$.
}

\begin{dimostrazione}
\passo{$g$ crescente ($\dot g > 0$): la CDF} Se $g$ è crescente, $g(X) \leq y$ succede esattamente quando $X \leq
g^{-1}(y)$:
\[ F_Y(y) = P(Y \leq y) = P(g(X) \leq y) = P\big(X \leq g^{-1}(y)\big) = F_X\big(g^{-1}(y)\big) = F_X\big(x(y)\big) . \]
\passo{$g$ crescente: la densità} Si deriva con la regola della catena e la derivata dell'inversa:
\[ f_Y(y) = \frac{\dd}{\dd y} F_Y(y) = f_X\big(g^{-1}(y)\big)\, \frac{\dd}{\dd y} g^{-1}(y) = \frac{f_X\big(g^{-1}(y)\big)}{\dot
   g\big(g^{-1}(y)\big)} . \]
\passo{$g$ decrescente ($\dot g < 0$): la CDF} Ora la disuguaglianza si rovescia, perché una funzione decrescente
scambia l'ordine: $g(X) \leq y$ succede esattamente quando $X \geq g^{-1}(y)$:
\[ F_Y(y) = P(g(X) \leq y) = P\big(X \geq g^{-1}(y)\big) = 1 - F_X\big(g^{-1}(y)\big) . \]
\passo{$g$ decrescente: la densità}
\[ f_Y(y) = -f_X\big(g^{-1}(y)\big)\, \frac{\dd}{\dd y} g^{-1}(y) = -\frac{f_X\big(g^{-1}(y)\big)}{\dot g\big(g^{-1}(y)\big)} . \]
\passo{Una formula sola} Nel secondo caso $\dot g < 0$, quindi $-\dot g = \abs{\dot g}$; nel primo $\dot g = \abs{\dot g}$.
In entrambi i casi il denominatore è $\abs{\dot g}$.
\end{dimostrazione}

\begin{figure}[H]
\centering
\begin{tikzpicture}
\begin{axis}[quaderno, width=8.4cm, height=6.6cm, xmin=0, xmax=3.3, ymin=0, ymax=3.1, xtick={1,2.4},
  xticklabels={$x_1$,$x_2$}, ytick={0.65,2.27}, yticklabels={$y_1$,$y_2$}, xlabel=$x$, ylabel=$y$, grid=none]
  \addplot[mspiA, domain=0.1:3.1, samples=60] {0.35*x^2+0.3};
  \fill[mspiB!35] (axis cs:0.95,0) rectangle (axis cs:1.05,0.65);
  \fill[mspiB!35] (axis cs:0,0.615) rectangle (axis cs:1,0.685);
  \fill[mspiC!35] (axis cs:2.35,0) rectangle (axis cs:2.45,2.27);
  \fill[mspiC!35] (axis cs:0,2.19) rectangle (axis cs:2.4,2.35);
  \node[font=\scriptsize, align=left, anchor=west] at (axis cs:1.15,0.35) {qui la curva è piatta:\\ $\mathrm{d}y$ stretto,
    densità alta};
  \node[font=\scriptsize, align=left, anchor=west] at (axis cs:0.1,2.65) {qui è ripida: $\mathrm{d}y$ largo,\\ densità bassa};
\end{axis}
\end{tikzpicture}
\caption{Perché si divide per $\lvert \dot g \rvert$. La probabilità contenuta in una striscia larga $\mathrm{d}x$ attorno a $x$,
cioè $f_X(x)\,\mathrm{d}x$, finisce in una striscia larga $\mathrm{d}y = \lvert \dot g(x) \rvert\,\mathrm{d}x$ attorno a
$y = g(x)$. La probabilità si conserva, $f_Y(y)\,\mathrm{d}y = f_X(x)\,\mathrm{d}x$, quindi
$f_Y = f_X/\lvert \dot g \rvert$: dove $g$ è ripida la probabilità si «spalma» su un intervallo più largo e la densità si
abbassa.}
\end{figure}

\begin{esempio}{Una trasformazione lineare}
Sia $Y = aX + b$ con $a \neq 0$. Allora $g(x) = ax + b$, $x(y) = \frac{y - b}{a}$ e $\dot g = a$:
\[ f_Y(y) = \frac{1}{\abs{a}}\, f_X\left( \frac{y - b}{a} \right) . \]
Il grafico di $f_Y$ è quello di $f_X$ spostato di $b$, allargato di un fattore $\abs{a}$ in orizzontale e abbassato
dello stesso fattore, così che l'area resti 1. Per esempio, se $X \sim \mathcal{C}(0, 1)$, allora
$bX + a \sim \mathcal{C}(a, b)$: i parametri della Cauchy sono proprio una traslazione e un cambio di scala. Useremo questa
formula nel Capitolo~15 per costruire tutte le gaussiane a partire da una sola.
\end{esempio}

\begin{esempio}{Il faro e la variabile di Cauchy}
Un faro ruota e si ferma puntando in una direzione a caso, con un angolo $\Theta \sim \mathcal{U}\left(-\frac{\pi}{2},
\frac{\pi}{2}\right)$ rispetto alla perpendicolare a un muro lungo e diritto che si trova a distanza 1. Il punto
illuminato sul muro è $X = \tan\Theta$. La funzione $g(\theta) = \tan\theta$ è crescente su
$\left(-\frac{\pi}{2}, \frac{\pi}{2}\right)$, con inversa $\theta(x) = \arctan x$ e derivata
$\dot g = 1 + \tan^2\theta = 1 + x^2$ (Capitolo~11). La densità di $\Theta$ è $\frac1\pi$ sull'intervallo, quindi
\[ f_X(x) = \frac{f_\Theta(\arctan x)}{1 + x^2} = \frac{1}{\pi}\, \frac{1}{1 + x^2} : \]
il punto illuminato è una variabile di Cauchy $\mathcal{C}(0, 1)$. Quando il faro punta quasi parallelo al muro, il
punto illuminato schizza lontanissimo: sono i valori enormi che impediscono alla Cauchy di avere una media.
\begin{center}
\begin{tikzpicture}
  \draw[line width=1.5pt, mspigray] (-3.6,2) -- (3.6,2);
  \node[font=\scriptsize, mspigray, anchor=south] at (-3.2,2) {muro};
  \fill[mspiB] (0,0) circle (0.12); \node[font=\scriptsize, below] at (0,-0.12) {faro};
  \draw[mspigray, dashed] (0,0) -- (0,2);
  \draw[mspiB, line width=0.8pt] (0,0) -- (1.6,2); \fill[evid!70!black] (1.6,2) circle (0.07);
  \draw[mspiB] (0,0.6) arc (90:51.3:0.6); \node[font=\scriptsize] at (0.3,0.85) {$\Theta$};
  \draw[decorate, decoration={brace, amplitude=4pt}] (0,2.1) -- (1.6,2.1) node[midway, above=4pt, font=\scriptsize]
    {$X = \tan\Theta$};
  \node[font=\scriptsize, anchor=east] at (-0.05,1.3) {1};
\end{tikzpicture}
\end{center}
\end{esempio}

\subsection{Secondo caso: \texorpdfstring{$g$}{g} non è invertibile}

Se $g$ non è invertibile, per un valore $y$ l'equazione $g(x) = y$ può avere più soluzioni:
\[ \forall y \in \mathcal{Y} \quad \exists\ \{x_i(y)\}_{i=1}^{M(y)} : \quad g[x_i(y)] = y . \]
Per fissare le idee supponiamo $M(y) = 4$, come nella figura, con $\dot g(x_1) < 0$. Allora $g(x) \leq y$ per
$x_1(y) \leq x \leq x_2(y)$; ovviamente sarà $\dot g(x_2) > 0$ e $g(x) \geq y$ per $x_2(y) \leq x \leq x_3(y)$. La
funzione ripassa per il valore $g(x_3) = y$ (con derivata negativa) e si mantiene al di sotto di $y$ fino a $x_4(y)$.

\begin{figure}[H]
\centering
\begin{tikzpicture}
\begin{axis}[quaderno, width=10cm, height=5.6cm, xmin=0, xmax=4.6, ymin=0, ymax=2.5, xtick={0.5,1.5,3,4},
  xticklabels={$x_1$,$x_2$,$x_3$,$x_4$}, ytick={1}, yticklabels={$y$}, xlabel=$x$, grid=none]
  \fill[evid] (axis cs:0.5,0) rectangle (axis cs:1.5,0.07);
  \fill[evid] (axis cs:3,0) rectangle (axis cs:4,0.07);
  \draw[mspiB, dashed] (axis cs:0,1) -- (axis cs:4.5,1);
  \addplot[mspiA, domain=0.2:4.3, samples=120] {1+0.3*(x-0.5)*(x-1.5)*(x-3)*(x-4)};
  \foreach \x in {0.5,1.5,3,4} {
    \edef\tmp{\noexpand\draw[mspigray, dotted] (axis cs:\x,0) -- (axis cs:\x,1);}\tmp
  }
  \addplot[only marks, mark=*, mark size=1.8pt, mspiB] coordinates {(0.5,1) (1.5,1) (3,1) (4,1)};
  \node[font=\scriptsize, anchor=south] at (axis cs:1,0.08) {$g(x) \leq y$};
  \node[font=\scriptsize, anchor=south] at (axis cs:3.5,0.08) {$g(x) \leq y$};
  \node[font=\small, mspiA!85!black] at (axis cs:2.25,1.75) {$g(x)$};
\end{axis}
\end{tikzpicture}
\caption{Una funzione non invertibile: il livello $y$ viene raggiunto in quattro punti. L'evento $\{Y \leq y\}$ è
l'unione dei due intervalli gialli dell'asse $x$.}
\end{figure}

Quindi
\[ F_Y(y) = P(Y \leq y) = P\big( \{x_1(y) \leq X \leq x_2(y)\} \cup \{x_3(y) \leq X \leq x_4(y)\} \big) , \]
e sfruttando il fatto che gli intervalli sono disgiunti
\[ F_Y(y) = F_X\big(x_2(y)\big) - F_X\big(x_1(y)\big) + F_X\big(x_4(y)\big) - F_X\big(x_3(y)\big) . \]

\teorema[nonmonotona]{Densità di una funzione non invertibile}{
  Se $y = g(x)$ ha le soluzioni $x_1(y), \ldots, x_M(y)$, tutte con $\dot g \neq 0$:
  \[ f_Y(y) = \sum_{i=1}^{M(y)} \frac{f_X\big(x_i(y)\big)}{\abs{\dot g\big(x_i(y)\big)}} . \]
}

\begin{dimostrazione}[nel caso della figura]
\passo{Si deriva la CDF} Ogni soluzione $x_i(y)$ è, vicino a $y$, l'inversa del tratto monotono di $g$ che passa per
$x_i$; per la derivata dell'inversa $\dot x_i(y) = \frac{1}{\dot g(x_i(y))}$. Derivando la CDF termine per termine con la
regola della catena:
\[ f_Y(y) = \frac{f_X(x_2)}{\dot g(x_2)} - \frac{f_X(x_1)}{\dot g(x_1)} + \frac{f_X(x_4)}{\dot g(x_4)} - \frac{f_X(x_3)}{\dot
   g(x_3)} . \]
\passo{I segni} Nel secondo e nel quarto termine la derivata a denominatore è negativa ($g$ scende in $x_1$ e in
$x_3$), quindi $-\frac{1}{\dot g} = \frac{1}{\abs{\dot g}}$; nel primo e nel terzo è positiva. Tutti i termini diventano
$\frac{f_X(x_i)}{\abs{\dot g(x_i)}}$.
\end{dimostrazione}

\begin{formula}[densità di una funzione di variabile aleatoria]
\[ f_Y(y) = \sum_{i \,:\, g(x_i) = y} \frac{f_X(x_i)}{\abs{\dot g(x_i)}} \qquad \text{(si somma su tutte le soluzioni di } y = g(x)) \]
\end{formula}

È lo stesso principio del discreto, dove le probabilità dei punti con la stessa immagine si sommavano; in più, ogni
contributo è diviso per la «velocità» $\abs{\dot g}$ con cui $g$ attraversa il livello $y$.

\begin{esempio}{Il quadrato di due variabili diverse}
Sia $g(x) = x^2$, e si considerino $X_1 \sim \mathcal{E}(\lambda)$ e $X_2 \sim \mathcal{C}(0, 1)$.

\textbf{Esponenziale.} Nel caso $Y_1 = X_1^2$, la variabile $X_1$ è non negativa ($\mathcal{X}_1 = [0, \infty[$), e su
$x \geq 0$ la trasformazione $g(x) = x^2$ è biunivoca: l'equazione $x^2 = y$ ammette solo la soluzione $x(y) = \sqrt{y}$.
Quindi
\[ \dot x(y) = \frac{1}{\dot g(x)} \bigg|_{x = \sqrt{y}} = \frac{1}{2x} \bigg|_{x = \sqrt{y}} = \frac{1}{2\sqrt{y}} ,
   \qquad f_{Y_1}(y) = \lambda e^{-\lambda x} u(x) \Big|_{x = \sqrt{y}}\, \dot x(y) = \frac{\lambda e^{-\lambda
   \sqrt{y}}}{2\sqrt{y}}\, u(y) . \]
\textbf{Cauchy.} Nel caso $Y_2 = X_2^2$ invece $\mathcal{X}_2 = \R$, la trasformazione non è biunivoca e per $y > 0$
l'equazione $x^2 = y$ ha due soluzioni, $x_{1,2}(y) = \pm\sqrt{y}$, in cui $\abs{\dot g} = 2\sqrt{y}$. Con
$f_{X_2}(x) = \frac{1}{\pi} \frac{1}{1 + x^2}$, che vale lo stesso in $\pm\sqrt{y}$:
\[ f_{Y_2}(y) = \left[ \frac{f_{X_2}(\sqrt{y})}{2\sqrt{y}} + \frac{f_{X_2}(-\sqrt{y})}{2\sqrt{y}} \right] u(y) =
   \frac{1}{\pi}\, \frac{1}{2\sqrt{y}}\, \frac{1}{1 + y} \cdot 2\; u(y) = \frac{1}{\pi \sqrt{y}\,(1 + y)}\, u(y) . \]
\begin{center}
\begin{tikzpicture}
\begin{axis}[quaderno, width=10cm, height=5cm, xmin=0, xmax=4.4, ymin=0, ymax=1.6, xtick={1,2,3,4}, ytick={0.5,1},
  yticklabels={$0{,}5$,$1$}, xlabel=$y$, legend style={font=\scriptsize, at={(0.98,0.98)}, anchor=north east}]
  \addplot[mspiA, domain=0.04:4.2, samples=120] {exp(-sqrt(x))/(2*sqrt(x))};
  \addplot[mspiB, domain=0.04:4.2, samples=120] {1/(pi*sqrt(x)*(1+x))};
  \legend{$f_{Y_1}$ ($\lambda = 1$), $f_{Y_2}$}
\end{axis}
\end{tikzpicture}
\end{center}
\end{esempio}

\warning{Una densità può valere infinito}{
  Entrambe le densità dell'esempio tendono a infinito per $y \to 0$, a causa del fattore $\frac{1}{\sqrt{y}}$: vicino
  a $x = 0$ la parabola $x^2$ è piattissima ($\dot g = 2x \to 0$), e schiaccia tanta probabilità in un intervallino di
  $y$. Non c'è nessuna contraddizione: l'area resta finita, perché $\int_0^1 \frac{1}{\sqrt{y}} \dd y = \left[ 2\sqrt{y}
  \right]_0^1 = 2$. Ancora una volta: conta l'area, non l'altezza.
}

\subsection{Due trasformazioni speciali}

\begin{esempio}{Ogni variabile, passata nella propria CDF, diventa uniforme}
Sia $X \sim f_X(x)$, e consideriamo $Y = F_X(X)$, cioè $g(x) = F_X(x)$. Qualunque sia la densità, $Y \in [0, 1]$,
perché la CDF assume valori fra 0 e 1. Inoltre $g$ è crescente; potrebbe non essere strettamente crescente (avere
tratti in cui è costante, dove $f_X = 0$), ma escludiamo questo caso. Allora
\[ g(x) = F_X(x) \;\Longrightarrow\; x(y) = F_X^{-1}(y) , \qquad \dot g(x) = f_X(x) , \]
e con la formula del primo caso
\[ f_Y(y) = \frac{f_X\big(F_X^{-1}(y)\big)}{f_X\big(F_X^{-1}(y)\big)} = 1 , \quad y \in [0, 1] , \qquad\text{cioè}\quad
   f_Y(y) = \mathrm{rect}\left( y - \tfrac12 \right) : \]
$Y$ è una variabile aleatoria uniforme in $(0, 1)$, $Y \sim \mathcal{U}(0, 1)$. Il motivo intuitivo: $F_X(X)$ è «la
frazione di probabilità che sta a sinistra del valore osservato», e questa frazione può essere un numero qualsiasi fra 0
e 1 con uguale probabilità.
\end{esempio}

Il risultato si può usare al contrario, ed è così che si simula al computer una variabile con qualsiasi densità.

\teorema[inversione]{Il metodo dell'inversione}{
  Se $U \sim \mathcal{U}(0, 1)$ e $F_X$ è una CDF strettamente crescente, la variabile $Z = F_X^{-1}(U)$ ha CDF $F_X$.
}

\begin{dimostrazione}
\passo{} Per ogni $z$:
\[ F_Z(z) = P(Z \leq z) = P\big(F_X^{-1}(U) \leq z\big) \underset{(a)}{=} P\big(U \leq F_X(z)\big) = F_U\big(F_X(z)\big)
   \underset{(b)}{=} F_X(z) . \]
\passo{(a)} $F_X$ è per ipotesi strettamente crescente, e quindi lo è anche la sua inversa: applicando $F_X$ ai due
membri di $F_X^{-1}(U) \leq z$ la disuguaglianza non cambia verso.
\passo{(b)} La CDF di una variabile uniforme in $(0, 1)$ è la funzione identica nell'intervallo $(0, 1)$:
$U \sim \mathcal{U}(0, 1) \iff f_U(u) = \mathrm{rect}\left(u - \frac12\right) \iff F_U(u) = u$ per $0 \leq u \leq 1$.
Quindi $F_U(F_X(z)) = F_X(z)$.
\end{dimostrazione}

I generatori di numeri casuali dei computer producono numeri uniformi in $(0, 1)$; con il metodo dell'inversione si
trasformano in numeri con la densità voluta. Questo ha notevoli conseguenze nelle procedure di simulazione dei sistemi
numerici (tutte le simulazioni di questa dispensa sono state fatte così).

\begin{esempio}{Generare tempi di attesa esponenziali}
Per $X \sim \mathcal{E}(\lambda)$, $F_X(x) = 1 - e^{-\lambda x}$. Risolvendo $u = 1 - e^{-\lambda x}$ rispetto a $x$:
$e^{-\lambda x} = 1 - u$, cioè $x = -\frac{1}{\lambda} \ln(1 - u)$. Quindi
\[ Z = -\frac{1}{\lambda} \ln(1 - U) \sim \mathcal{E}(\lambda) . \]
\begin{center}
\begin{tikzpicture}
\begin{groupplot}[group style={group size=2 by 1, horizontal sep=1.6cm},
  quaderno, width=6.8cm, height=4.8cm, xmin=0, xmax=5.3, xtick={1,2,3,4,5}]
\nextgroupplot[ymin=0, ymax=1.15, ytick={0.7,1}, yticklabels={$u$,$1$}, xlabel=$x$, ylabel={$F_X(x)$}]
  \addplot[mspiB, domain=0:5.1, samples=80] {1-exp(-x)};
  \draw[-{Stealth}, mspiA, line width=0.8pt] (axis cs:0,0.7) -- (axis cs:1.204,0.7);
  \draw[-{Stealth}, mspiA, line width=0.8pt] (axis cs:1.204,0.7) -- (axis cs:1.204,0);
  \node[font=\scriptsize, mspiA!85!black, anchor=north west] at (axis cs:1.25,0.12) {$z = F_X^{-1}(u)$};
\nextgroupplot[ymin=0, ymax=1.1, ytick={0.5,1}, yticklabels={$0{,}5$,$1$}, xlabel=$x$]
  \addplot[ybar interval, fill=mspiA!35, draw=mspiA!70] coordinates {(0.00,0.8496) (0.25,0.6960) (0.50,0.5480) (0.75,0.4080) (1.00,0.3064) (1.25,0.2808) (1.50,0.2040) (1.75,0.1536) (2.00,0.1136) (2.25,0.1024) (2.50,0.0648) (2.75,0.0664) (3.00,0.0512) (3.25,0.0280) (3.50,0.0352) (3.75,0.0184) (4.00,0.0112) (4.25,0.0136) (4.50,0.0088) (4.75,0.0080) (5.00,0.0080)};
  \addplot[mspiB, line width=1.1pt, domain=0:5.1, samples=80] {exp(-x)};
\end{groupplot}
\end{tikzpicture}
\end{center}
A sinistra: si estrae $u$ uniforme sull'asse verticale, si va orizzontalmente fino alla CDF e si scende: il punto
d'arrivo è $z$. A destra: l'istogramma di 5000 numeri ottenuti così (con $\lambda = 1$) segue la densità esponenziale.
Dove la CDF è ripida (vicino a 0) molti valori di $u$ finiscono in un intervallo stretto di $x$: lì la densità è alta.
\end{esempio}

\begin{esempio}{La sinusoide con fase a caso}
Siano $X \sim \mathcal{U}\left(-\frac12, \frac12\right)$ e $g(x) = A \cos(2\pi x + \varphi)$, con $A > 0$. Cerchiamo la
densità di $Y = A \cos(2\pi X + \varphi)$: il valore di una sinusoide osservata in un istante a caso.

Si noti che $\mathcal{Y} = [-A, A]$ (il coseno sta fra $-1$ e 1) e che la trasformazione non è biunivoca. Poiché
$\mathcal{X} = \left[-\frac12, \frac12\right]$, l'argomento del coseno varia in un intervallo di ampiezza $2\pi$, un
periodo intero, e in questo intervallo l'equazione $y = A \cos(2\pi x + \varphi)$ ha due soluzioni $x_1, x_2$ (una
mentre il coseno scende e una mentre sale). Quindi
\[ f_Y(y) = \frac{f_X(x_1(y))}{\abs{\dot g(x_1(y))}} + \frac{f_X(x_2(y))}{\abs{\dot g(x_2(y))}} \underset{(a)}{=}
   \frac{1}{\abs{\dot g(x_1(y))}} + \frac{1}{\abs{\dot g(x_2(y))}} , \]
dove (a) vale perché $f_X = 1$ sull'intervallo. La derivata del coseno è meno il seno (con la regola della catena il
fattore $2\pi$ esce dall'argomento):
\[ \dot g(x) = -A \sin(2\pi x + \varphi) \cdot 2\pi , \qquad \abs{\dot g(x)} = 2\pi A \abs{\sin(2\pi x + \varphi)} . \]
Va scritta in funzione di $y$: per Pitagora $\sin^2\alpha + \cos^2\alpha = 1$ (seno e coseno sono i cateti di un
triangolo rettangolo con ipotenusa 1), quindi $\abs{\sin\alpha} = \sqrt{1 - \cos^2\alpha}$, e $\cos(2\pi x + \varphi) =
\frac{y}{A}$ nei punti $x_i(y)$:
\[ \abs{\dot g(x_i(y))} = 2\pi A \sqrt{1 - \cos^2(2\pi x_i + \varphi)} = 2\pi A \sqrt{1 - \left( \frac{y}{A} \right)^2} . \]
I due termini sono uguali, e
\[ f_Y(y) = \frac{2}{2\pi A \sqrt{1 - \left( \frac{y}{A} \right)^2}} = \frac{1}{\pi A \sqrt{1 - \left( \frac{y}{A} \right)^2}} =
   \frac{1}{\pi \sqrt{A^2 - y^2}} , \qquad y \in (-A, A) . \]
È la \term{distribuzione dell'arcoseno}, a forma di «U»: la sinusoide passa più tempo vicino ai massimi e ai minimi, dove
rallenta per invertire il moto, che vicino allo zero, dove corre.
\begin{center}
\begin{tikzpicture}
\begin{groupplot}[group style={group size=2 by 1, horizontal sep=1.4cm},
  quaderno, width=6.8cm, height=4.6cm]
\nextgroupplot[xmin=0, xmax=1.1, ymin=-1.35, ymax=1.35, xtick={0.5,1}, ytick={-1,1}, yticklabels={$-A$,$A$}, xlabel=$t$]
  \fill[evid] (axis cs:0,0.85) rectangle (axis cs:1,1);
  \fill[mspiC!25] (axis cs:0,-0.075) rectangle (axis cs:1,0.075);
  \addplot[mspiA, domain=0:1, samples=100] {cos(deg(2*pi*x))};
  \addplot[mspiA, line width=2.2pt, domain=0:0.0877, samples=10] {cos(deg(2*pi*x))};
  \addplot[mspiA, line width=2.2pt, domain=0.9123:1, samples=10] {cos(deg(2*pi*x))};
  \addplot[mspiA, line width=2.2pt, domain=0.238:0.262, samples=5] {cos(deg(2*pi*x))};
  \addplot[mspiA, line width=2.2pt, domain=0.738:0.762, samples=5] {cos(deg(2*pi*x))};
\nextgroupplot[xmin=-1.15, xmax=1.15, ymin=0, ymax=1.6, xtick={-1,1}, xticklabels={$-A$,$A$}, ytick={0.3183},
  yticklabels={$\frac{1}{\pi A}$}, xlabel=$y$]
  \addplot[ybar interval, fill=mspiA!30, draw=mspiA!60] coordinates {(-1.00,1.4325) (-0.90,0.6020) (-0.80,0.4835) (-0.70,0.4310) (-0.60,0.3805) (-0.50,0.3725) (-0.40,0.3460) (-0.30,0.3245) (-0.20,0.3160) (-0.10,0.3160) (0.00,0.3165) (0.10,0.3255) (0.20,0.3510) (0.30,0.3430) (0.40,0.3585) (0.50,0.3990) (0.60,0.4070) (0.70,0.4805) (0.80,0.5995) (0.90,1.4150) (1.00,1.4150)};
  \addplot[mspiB, line width=1.1pt, domain=-0.985:0.985, samples=120] {1/(pi*sqrt(1-x^2))};
\end{groupplot}
\end{tikzpicture}
\end{center}
A sinistra: in un periodo, la sinusoide resta nella fascia gialla vicino al massimo (le parti spesse a sinistra e a
destra) molto più a lungo che nella fascia verde, alta uguale, attorno allo zero (le due parti spesse in mezzo). A
destra: l'istogramma di 20\,000 valori simulati con $A = 1$ e la densità dell'arcoseno.
\end{esempio}

% ============================================================
\section{Terzo caso: la conversione A/D}
% ============================================================

Quando $X$ è continua e $Y$ è discreta si ha una \textbf{conversione analogico-digitale} (A/D) della variabile
aleatoria. Supponiamo di voler rappresentare $X$ (con perdite) con $R$ bit, cioè con $M = 2^R$ livelli.
\begin{itemize}
  \item Si scelgono gli $M$ valori di rappresentazione $\mathcal{Y} = \{y_1, y_2, \ldots, y_M\}$, i valori che $Y$ può
        assumere.
  \item Si divide l'alfabeto $\mathcal{X}$ in $M$ intervalli, definiti dalle $M + 1$ soglie $\{x_i\}_{i=1}^{M+1}$.
\end{itemize}

\dfn{Rappresentazione a R bit}{
  Si definisce \term{rappresentazione a $R$ bit} di $X$ la variabile aleatoria discreta
  \[ Y = y_i \qquad \text{se} \quad x_i \leq X < x_{i+1} , \qquad i = 1, 2, \ldots, M . \]
}

La pmf di $Y$ si scrive facilmente: $Y$ vale $y_i$ quando $X$ cade nell'$i$-esimo intervallo, quindi
\[ p_Y(y_i) = P(x_i \leq X < x_{i+1}) = F_X(x_{i+1}) - F_X(x_i) , \qquad i = 1, 2, \ldots, M . \]
Ovviamente tanto la partizione quanto i livelli di rappresentazione sono gradi di libertà a disposizione del
progettista: scegliendoli bene si riduce l'errore commesso rappresentando $X$ con $Y$.

\begin{esempio}{Un convertitore a 3 bit}
I campioni di un segnale vocale si modellano bene con una variabile di Laplace (valori piccoli frequentissimi, valori
grandi rari). Sia $X \sim \mathcal{L}(1)$, quantizzata con $R = 3$ bit, cioè $M = 8$ livelli: soglie
$-\infty, -3, -2, -1, 0, 1, 2, 3, +\infty$ e livelli $-3{,}5;\ -2{,}5;\ \ldots;\ 3{,}5$ (il centro di ogni intervallo,
e per i due intervalli estremi il punto a distanza $\frac12$ dall'ultima soglia).
\begin{center}
\begin{tikzpicture}
\begin{groupplot}[group style={group size=2 by 1, horizontal sep=1.5cm}, quaderno, axis lines=left,
  width=6.8cm, height=5.2cm]
\nextgroupplot[xmin=-4.6, xmax=4.6, ymin=-4, ymax=4, xtick={-3,-2,-1,0,1,2,3}, ytick={-3.5,-1.5,0.5,1.5,3.5},
  yticklabels={$-3{,}5$,$-1{,}5$,$0{,}5$,$1{,}5$,$3{,}5$}, xlabel=$x$, ylabel={$y = g(x)$}]
  \addplot[mspiA, line width=1.3pt, jump mark left] coordinates {(-4.5,-3.5) (-3,-2.5) (-2,-1.5) (-1,-0.5) (0,0.5) (1,1.5) (2,2.5) (3,3.5) (4.5,3.5)};
  \addplot[mspiB, domain=-4.5:4.5, samples=121] {3*0.5*exp(-abs(x))-3.9};
  \node[font=\scriptsize, mspiB!85!black, anchor=west] at (axis cs:1.2,-3.1) {$f_X(x)$};
\nextgroupplot[xmin=-4.2, xmax=4.2, ymin=0, ymax=0.38, xtick={-3.5,-2.5,-1.5,-0.5,0.5,1.5,2.5,3.5},
  xticklabels={$-3{,}5$,,$-1{,}5$,,,$1{,}5$,,$3{,}5$}, ytick={0.1,0.2,0.3}, yticklabels={$0{,}1$,$0{,}2$,$0{,}3$},
  xlabel=$y$, ylabel={$p_Y(y)$}]
  \addplot[bastoncini, mspiA] coordinates {(-3.5,0.0249) (-2.5,0.0428) (-1.5,0.1163) (-0.5,0.3161) (0.5,0.3161) (1.5,0.1163) (2.5,0.0428) (3.5,0.0249)};
\end{groupplot}
\end{tikzpicture}
\end{center}
A sinistra la caratteristica «a scalini» del convertitore (in basso, per riferimento, la densità di $X$); a destra la
pmf di $Y$. Per esempio
\[ p_Y(0{,}5) = F_X(1) - F_X(0) = \left( 1 - \tfrac12 e^{-1} \right) - \tfrac12 \approx 0{,}316 , \qquad
   p_Y(3{,}5) = F_X(+\infty) - F_X(3) = \tfrac12 e^{-3} \approx 0{,}025 . \]
Più del 63\% delle volte il convertitore usa uno dei due livelli centrali, e solo il 5\% delle volte uno dei due
estremi. Nella parte del corso sull'informazione vedremo che questa disuguaglianza si può sfruttare per comprimere i
dati.
\end{esempio}

\gbox{In sintesi}{mspiA!80!black}{
\begin{itemize}
  \item \textbf{Condizionamento}: $F_{X \mid A}(x) = P(\{X \leq x\} \cap A)/P(A)$ e $f_{X \mid A} = F_{X \mid A}'$; se $A$
        è un intervallo, la densità viene tagliata e rialzata dividendo per $P(A)$.
  \item \textbf{Probabilità totale}: $f_X = \sum_m f_{X \mid E_m} P(E_m)$ e $E[X] = \sum_m E[X \mid E_m] P(E_m)$.
  \item Se $Y = g(X)$: $f_Y(y) = \sum_i \frac{f_X(x_i)}{\abs{\dot g(x_i)}}$, sommando sulle soluzioni di $g(x) = y$. La
        probabilità si conserva, $f_Y \mathrm{d}y = f_X \mathrm{d}x$: dove $g$ è ripida la densità si abbassa. Lineare:
        $f_{aX + b}(y) = f_X\left(\frac{y - b}{a}\right)/\abs{a}$.
  \item $F_X(X) \sim \mathcal{U}(0, 1)$, e viceversa $F_X^{-1}(U)$ ha CDF $F_X$: è il metodo dell'inversione per
        simulare qualsiasi densità.
  \item $A\cos(2\pi X + \varphi)$ con $X$ uniforme su un periodo ha la densità dell'arcoseno $\frac{1}{\pi\sqrt{A^2 - y^2}}$.
  \item \textbf{Conversione A/D} a $R$ bit: $M = 2^R$ livelli, $p_Y(y_i) = F_X(x_{i+1}) - F_X(x_i)$.
\end{itemize}
}
